\documentclass[11pt,a4paper]{article}
\usepackage{amsmath,amssymb,amsfonts,amsthm}
\usepackage{geometry}
\geometry{margin=1in}
\title{\textbf{Geotemporal Hydrodynamics (GTH) v12.0\\Paper 1: Density Ceiling and Core Regularization}}
\author{CoderQuan2}
\date{\today}
\begin{document}
\maketitle
\begin{abstract}
We establish the mathematical framework for core regularization within Geotemporal Hydrodynamics (GTH). By enforcing a strict density ceiling boundary condition $\rho \le \rho_{\max}$, gravitational and fluid-dynamic singularities are eliminated across non-linear manifold collapse regimes.
\end{abstract}
\section{Field Equations and Density Ceiling}
The modified continuous energy-momentum conservation tensor under GTH formulation takes the form:
\begin{equation}
T_{\mu\nu} = (\rho + P) u_\mu u_\nu + P g_{\mu\nu} + \Xi_{\mu\nu}(\rho/\rho_{\max})
\end{equation}
where $\Xi_{\mu\nu}$ represents the non-linear core regularization factor preventing singular field concentration as $\rho \to \rho_{\max}$.
\section{Core Boundary Conditions}
At $r \to 0$, metric components regularize continuously without horizon collapse:
\begin{equation}
ds^2 = - \left(1 - \frac{2M r^2}{(r^2 + r_c^2)^{3/2}}\right) dt^2 + \left(1 - \frac{2M r^2}{(r^2 + r_c^2)^{3/2}}\right)^{-1} dr^2 + r^2 d\Omega^2
\end{equation}
where $r_c = \sqrt{\frac{3M}{4\pi \rho_{\max}}}$.
\end{document}
